% Options for packages loaded elsewhere
% Options for packages loaded elsewhere
\PassOptionsToPackage{unicode}{hyperref}
\PassOptionsToPackage{hyphens}{url}
\PassOptionsToPackage{dvipsnames,svgnames,x11names}{xcolor}
%
\documentclass[
  french,
  9pt,
  a4paper,
  DIV=11,
  twoside=semi]{scrreprt}
\usepackage{xcolor}
\usepackage{amsmath,amssymb}
\setcounter{secnumdepth}{3}
\usepackage{iftex}
\ifPDFTeX
  \usepackage[T1]{fontenc}
  \usepackage[utf8]{inputenc}
  \usepackage{textcomp} % provide euro and other symbols
\else % if luatex or xetex
  \usepackage{unicode-math} % this also loads fontspec
  \defaultfontfeatures{Scale=MatchLowercase}
  \defaultfontfeatures[\rmfamily]{Ligatures=TeX,Scale=1}
\fi
\usepackage{lmodern}
\ifPDFTeX\else
  % xetex/luatex font selection
\fi
% Use upquote if available, for straight quotes in verbatim environments
\IfFileExists{upquote.sty}{\usepackage{upquote}}{}
\IfFileExists{microtype.sty}{% use microtype if available
  \usepackage[]{microtype}
  \UseMicrotypeSet[protrusion]{basicmath} % disable protrusion for tt fonts
}{}
\usepackage{setspace}
\makeatletter
\@ifundefined{KOMAClassName}{% if non-KOMA class
  \IfFileExists{parskip.sty}{%
    \usepackage{parskip}
  }{% else
    \setlength{\parindent}{0pt}
    \setlength{\parskip}{6pt plus 2pt minus 1pt}}
}{% if KOMA class
  \KOMAoptions{parskip=half}}
\makeatother
% Make \paragraph and \subparagraph free-standing
\makeatletter
\ifx\paragraph\undefined\else
  \let\oldparagraph\paragraph
  \renewcommand{\paragraph}{
    \@ifstar
      \xxxParagraphStar
      \xxxParagraphNoStar
  }
  \newcommand{\xxxParagraphStar}[1]{\oldparagraph*{#1}\mbox{}}
  \newcommand{\xxxParagraphNoStar}[1]{\oldparagraph{#1}\mbox{}}
\fi
\ifx\subparagraph\undefined\else
  \let\oldsubparagraph\subparagraph
  \renewcommand{\subparagraph}{
    \@ifstar
      \xxxSubParagraphStar
      \xxxSubParagraphNoStar
  }
  \newcommand{\xxxSubParagraphStar}[1]{\oldsubparagraph*{#1}\mbox{}}
  \newcommand{\xxxSubParagraphNoStar}[1]{\oldsubparagraph{#1}\mbox{}}
\fi
\makeatother


\usepackage{longtable,booktabs,array}
\newcounter{none} % for unnumbered tables
\usepackage{calc} % for calculating minipage widths
% Correct order of tables after \paragraph or \subparagraph
\usepackage{etoolbox}
\makeatletter
\patchcmd\longtable{\par}{\if@noskipsec\mbox{}\fi\par}{}{}
\makeatother
% Allow footnotes in longtable head/foot
\IfFileExists{footnotehyper.sty}{\usepackage{footnotehyper}}{\usepackage{footnote}}
\makesavenoteenv{longtable}
\usepackage{graphicx}
\makeatletter
\newsavebox\pandoc@box
\newcommand*\pandocbounded[1]{% scales image to fit in text height/width
  \sbox\pandoc@box{#1}%
  \Gscale@div\@tempa{\textheight}{\dimexpr\ht\pandoc@box+\dp\pandoc@box\relax}%
  \Gscale@div\@tempb{\linewidth}{\wd\pandoc@box}%
  \ifdim\@tempb\p@<\@tempa\p@\let\@tempa\@tempb\fi% select the smaller of both
  \ifdim\@tempa\p@<\p@\scalebox{\@tempa}{\usebox\pandoc@box}%
  \else\usebox{\pandoc@box}%
  \fi%
}
% Set default figure placement to htbp
\def\fps@figure{htbp}
\makeatother


% definitions for citeproc citations
\NewDocumentCommand\citeproctext{}{}
\NewDocumentCommand\citeproc{mm}{%
  \begingroup\def\citeproctext{#2}\cite{#1}\endgroup}
\makeatletter
 % allow citations to break across lines
 \let\@cite@ofmt\@firstofone
 % avoid brackets around text for \cite:
 \def\@biblabel#1{}
 \def\@cite#1#2{{#1\if@tempswa , #2\fi}}
\makeatother
\newlength{\cslhangindent}
\setlength{\cslhangindent}{1.5em}
\newlength{\csllabelwidth}
\setlength{\csllabelwidth}{3em}
\newenvironment{CSLReferences}[2] % #1 hanging-indent, #2 entry-spacing
 {\begin{list}{}{%
  \setlength{\itemindent}{0pt}
  \setlength{\leftmargin}{0pt}
  \setlength{\parsep}{0pt}
  % turn on hanging indent if param 1 is 1
  \ifodd #1
   \setlength{\leftmargin}{\cslhangindent}
   \setlength{\itemindent}{-1\cslhangindent}
  \fi
  % set entry spacing
  \setlength{\itemsep}{#2\baselineskip}}}
 {\end{list}}
\usepackage{calc}
\newcommand{\CSLBlock}[1]{\hfill\break\parbox[t]{\linewidth}{\strut\ignorespaces#1\strut}}
\newcommand{\CSLLeftMargin}[1]{\parbox[t]{\csllabelwidth}{\strut#1\strut}}
\newcommand{\CSLRightInline}[1]{\parbox[t]{\linewidth - \csllabelwidth}{\strut#1\strut}}
\newcommand{\CSLIndent}[1]{\hspace{\cslhangindent}#1}

\ifLuaTeX
\usepackage[bidi=basic,shorthands=off]{babel}
\else
\usepackage[bidi=default,shorthands=off]{babel}
\fi
\ifLuaTeX
  \usepackage{selnolig} % disable illegal ligatures
\fi


\setlength{\emergencystretch}{3em} % prevent overfull lines

\providecommand{\tightlist}{%
  \setlength{\itemsep}{0pt}\setlength{\parskip}{0pt}}



 


\usepackage{booktabs}
\usepackage{caption}
\usepackage{longtable}
\usepackage{colortbl}
\usepackage{array}
\usepackage{anyfontsize}
\usepackage{multirow}
\makeatletter
\@ifpackageloaded{caption}{}{\usepackage{caption}}
\AtBeginDocument{%
\ifdefined\contentsname
  \renewcommand*\contentsname{Table des matières}
\else
  \newcommand\contentsname{Table des matières}
\fi
\ifdefined\listfigurename
  \renewcommand*\listfigurename{Graphiques}
\else
  \newcommand\listfigurename{Graphiques}
\fi
\ifdefined\listtablename
  \renewcommand*\listtablename{Tableaux}
\else
  \newcommand\listtablename{Tableaux}
\fi
\ifdefined\figurename
  \renewcommand*\figurename{Graphique}
\else
  \newcommand\figurename{Graphique}
\fi
\ifdefined\tablename
  \renewcommand*\tablename{Tableau}
\else
  \newcommand\tablename{Tableau}
\fi
}
\@ifpackageloaded{float}{}{\usepackage{float}}
\floatstyle{ruled}
\@ifundefined{c@chapter}{\newfloat{codelisting}{h}{lop}}{\newfloat{codelisting}{h}{lop}[chapter]}
\floatname{codelisting}{Listing}
\newcommand*\listoflistings{\listof{codelisting}{Liste des Listings}}
\makeatother
\makeatletter
\makeatother
\makeatletter
\@ifpackageloaded{caption}{}{\usepackage{caption}}
\@ifpackageloaded{subcaption}{}{\usepackage{subcaption}}
\makeatother
\usepackage{bookmark}
\IfFileExists{xurl.sty}{\usepackage{xurl}}{} % add URL line breaks if available
\urlstyle{same}
\hypersetup{
  pdftitle={Ajustement budgétaire 2026-2030},
  pdfauthor={Xavier Timbeau},
  pdflang={fr},
  colorlinks=true,
  linkcolor={blue},
  filecolor={Maroon},
  citecolor={Blue},
  urlcolor={Blue},
  pdfcreator={LaTeX via pandoc}}


%%% title.tex we use to call other packages. Because this works
\usepackage{titlepic}
\usepackage{titling}
\usepackage{graphicx}
\usepackage{fontspec}
\usepackage{placeins}
\usepackage{graphbox}
\usepackage{tikz}
\usepackage{geometry}
\usepackage{xcolor}
\usepackage{amsmath}
\usepackage[some]{background}
\usepackage{lipsum}
\usepackage{caption}
\usepackage{datetime}
\usepackage{xstring}
\usepackage{blindtext}
\usepackage{scrlayer-scrpage}

\setmainfont{OpenSans}[
  UprightFont = {*-Regular},
  BoldFont = {*-Bold},
  BoldItalicFont = {*-BoldItalic},
  ItalicFont = {*-Italic},
  Path = {OpenSans/},
  Extension = {.ttf}
]
\setsansfont{OpenSans}[
  UprightFont = {*-Regular},
  BoldFont = {*-Bold},
  BoldItalicFont = {*-BoldItalic},
  ItalicFont = {*-Italic},
  Path = {OpenSans/},
  Extension = {.ttf}
]

\def\getYear#1{\StrLeft{#1}{4}}
\def\getMonth#1{\StrMid{#1}{6}{7}}
\def\getDay#1{\StrRight{#1}{2}}

\def\jolimois#1{\monthname{\getMonth{#1}}}

\definecolor{quarto-callout-color}{HTML}{eeeeee}
\definecolor{quarto-callout-tip-color}{HTML}{dddddd}
\definecolor{quarto-callout-tip-color-frame}{HTML}{eeeeee}

\clearpairofpagestyles

\KOMAoptions{headsepline=true}

\pagestyle{scrheadings}
\setkomafont{pageheadfoot}{\small}

\lofoot*{\thepage}
\refoot*{\thepage}
\begin{document}



\definecolor{ofcepbbleu}{RGB}{1, 97, 131}
\definecolor{ofcerouge}{RGB}{198, 45, 43}
\definecolor{scporouge}{RGB}{231, 0, 26}

\begin{titlepage}
  \backgroundsetup{
    scale=1,
    angle=0,
    opacity=1,
    contents={
  \begin{tikzpicture}[remember picture,overlay]
    \useasboundingbox (0,0) rectangle(\the\paperwidth,\the\paperheight);
      \node [anchor = center] at (-2.875-5.25,12.5) {\includegraphics[width=2.5cm]{../../\_extensions/ofce/note/logos/ofce\_m.png}};
      \node [anchor = center] at (-2.875-5.25,-12.75){\includegraphics[width=2.5cm]{../../\_extensions/ofce/note/logos/sciencespo.png}};
      \node [anchor = east] at (8.5,-11.5){\textcolor{black}{Publié le : 2026-07-06}};
            \end{tikzpicture}
  }
}
\BgThispage

\hspace{4cm}
\begin{minipage}{12.5cm}
  \vspace{5cm}
  \begin{flushleft}
    \begin{spacing}{2.5}
      \textcolor{scporouge}{\Huge\textbf{\textsf{Ajustement budgétaire
2026-2030}}}
    \end{spacing}
    \vspace{20mm}
  \end{flushleft}
  
    
   \textbf{Xavier Timbeau}
  % by-author
 \end{minipage}


\end{titlepage}
\renewcommand*\contentsname{Table des matières}
{
\setcounter{tocdepth}{0}
\tableofcontents
}

\setstretch{1.25}
\chapter{Principaux résultats}\label{principaux-ruxe9sultats}

L'ajustement «\,optimal\,»\,-- tel que défini plus loin, voir
principales hypothèses\,-- conduit à une impulsion de \textbf{1,2~points
de PIB} en 2027. C'est un résultat qui découle d'un déficit initial
important (5,1\% en 2025) et d'un objectif de dette à brève échéance
(2037). L'absence de réduction du déficit en 2026, justifié par le
contexte politique, accentue le report de l'effort et donc son
intensité. Les conséquences de ce fort ajustement sont le creusement de
l'écart de production en 2027 et en 2028. A partir de 2029, les
impulsions sont plus faibles et la dynamique du multiplicateur fait que
l'ajustement de 2027 perd en effet sur l'écart de production. Cependant,
parce qu'il y a chaque année une impulsion négative, même moindre,
l'écart de production reste négatif jusu'en 2035
(graphique~\ref{fig-optimal}).

\begin{table}

\caption{\label{tbl-optimal1}Ajustement optimal, impulsion nulle en
2026}

\centering{

\fontsize{6.0pt}{7.0pt}\selectfont
\begin{tabular*}{\linewidth}{@{\extracolsep{\fill}}>{\raggedright\arraybackslash}p{\dimexpr 165.00pt -2\tabcolsep-1.5\arrayrulewidth}>{\raggedleft\arraybackslash}p{\dimexpr 37.50pt -2\tabcolsep-1.5\arrayrulewidth}>{\raggedleft\arraybackslash}p{\dimexpr 37.50pt -2\tabcolsep-1.5\arrayrulewidth}>{\raggedleft\arraybackslash}p{\dimexpr 37.50pt -2\tabcolsep-1.5\arrayrulewidth}>{\raggedleft\arraybackslash}p{\dimexpr 37.50pt -2\tabcolsep-1.5\arrayrulewidth}>{\raggedleft\arraybackslash}p{\dimexpr 37.50pt -2\tabcolsep-1.5\arrayrulewidth}>{\raggedleft\arraybackslash}p{\dimexpr 37.50pt -2\tabcolsep-1.5\arrayrulewidth}>{\raggedleft\arraybackslash}p{\dimexpr 37.50pt -2\tabcolsep-1.5\arrayrulewidth}>{\raggedleft\arraybackslash}p{\dimexpr 37.50pt -2\tabcolsep-1.5\arrayrulewidth}>{\raggedleft\arraybackslash}p{\dimexpr 37.50pt -2\tabcolsep-1.5\arrayrulewidth}>{\raggedleft\arraybackslash}p{\dimexpr 37.50pt -2\tabcolsep-1.5\arrayrulewidth}}
\toprule
 & 2024 & 2025 & 2026 & 2027 & 2028 & 2029 & 2030 & 2031 & LT (2050) & Max. 25-50 \\ 
\midrule\addlinespace[2.5pt]
Croissance du PIB\textsuperscript{\textit{a}} &   1,19  &   0,84  & {\cellcolor[HTML]{F2F2F2}{  1,05 }} & {\cellcolor[HTML]{F2F2F2}{  0,04 }} & {\cellcolor[HTML]{F2F2F2}{  0,85 }} & {\cellcolor[HTML]{F2F2F2}{  1,17 }} & {\cellcolor[HTML]{F2F2F2}{  1,25 }} & {\cellcolor[HTML]{EBEBEB}{  1,25 }} & {\cellcolor[HTML]{EBEBEB}{  1,01 }} & {\cellcolor[HTML]{EBEBEB}{—}} \\ 
Ecart de production\textsuperscript{\textit{b}} &  -0,28  &   0,21  & {\cellcolor[HTML]{F2F2F2}{  0,08 }} & {\cellcolor[HTML]{F2F2F2}{ -1,03 }} & {\cellcolor[HTML]{F2F2F2}{ -1,33 }} & {\cellcolor[HTML]{F2F2F2}{ -1,21 }} & {\cellcolor[HTML]{F2F2F2}{ -0,99 }} & {\cellcolor[HTML]{EBEBEB}{ -0,78 }} & {\cellcolor[HTML]{EBEBEB}{  0,02 }} & \textsuperscript{\textit{c}}{\cellcolor[HTML]{EBEBEB}{ -1,33 }} \\ 
Croissance du PIB potentiel\textsuperscript{\textit{a}} & — & — & {\cellcolor[HTML]{F2F2F2}{  1,18 }} & {\cellcolor[HTML]{F2F2F2}{  1,16 }} & {\cellcolor[HTML]{F2F2F2}{  1,15 }} & {\cellcolor[HTML]{F2F2F2}{  1,04 }} & {\cellcolor[HTML]{F2F2F2}{  1,03 }} & {\cellcolor[HTML]{EBEBEB}{  1,03 }} & {\cellcolor[HTML]{EBEBEB}{  1,00 }} & {\cellcolor[HTML]{EBEBEB}{—}} \\ 
Inflation\textsuperscript{\textit{d}} &   2,09  &   1,18  & {\cellcolor[HTML]{F2F2F2}{  1,82 }} & {\cellcolor[HTML]{F2F2F2}{  1,78 }} & {\cellcolor[HTML]{F2F2F2}{  1,76 }} & {\cellcolor[HTML]{F2F2F2}{  1,61 }} & {\cellcolor[HTML]{F2F2F2}{  1,56 }} & {\cellcolor[HTML]{EBEBEB}{  1,52 }} & {\cellcolor[HTML]{EBEBEB}{  1,72 }} & {\cellcolor[HTML]{EBEBEB}{—}} \\ 
Taux de chômage\textsuperscript{\textit{e}} &   7,40  &   7,70  & {\cellcolor[HTML]{F2F2F2}{  7,78 }} & {\cellcolor[HTML]{F2F2F2}{  8,28 }} & {\cellcolor[HTML]{F2F2F2}{  8,47 }} & {\cellcolor[HTML]{F2F2F2}{  8,45 }} & {\cellcolor[HTML]{F2F2F2}{  8,26 }} & {\cellcolor[HTML]{EBEBEB}{  8,03 }} & {\cellcolor[HTML]{EBEBEB}{  6,98 }} & {\cellcolor[HTML]{EBEBEB}{  8,47 }} \\ 
{Solde public}\textsuperscript{\textit{b}} & { -5,93 } & { -5,09 } & {\cellcolor[HTML]{F2F2F2}{ -4,95 }} & {\cellcolor[HTML]{F2F2F2}{ -4,62 }} & {\cellcolor[HTML]{F2F2F2}{ -4,39 }} & {\cellcolor[HTML]{F2F2F2}{ -4,02 }} & {\cellcolor[HTML]{F2F2F2}{ -3,66 }} & {\cellcolor[HTML]{EBEBEB}{ -3,35 }} & {\cellcolor[HTML]{EBEBEB}{ -3,20 }} & {\cellcolor[HTML]{EBEBEB}{ -2,89 }} \\ 
Solde public primaire\textsuperscript{\textit{b}} &  -3,73  &  -2,75  & {\cellcolor[HTML]{F2F2F2}{ -2,58 }} & {\cellcolor[HTML]{F2F2F2}{ -1,98 }} & {\cellcolor[HTML]{F2F2F2}{ -1,55 }} & {\cellcolor[HTML]{F2F2F2}{ -1,05 }} & {\cellcolor[HTML]{F2F2F2}{ -0,57 }} & {\cellcolor[HTML]{EBEBEB}{ -0,16 }} & {\cellcolor[HTML]{EBEBEB}{  0,58 }} & {\cellcolor[HTML]{EBEBEB}{  0,59 }} \\ 
Solde public structurel primaire\textsuperscript{\textit{b}} &  -3,61  &  -2,85  & {\cellcolor[HTML]{F2F2F2}{ -2,62 }} & {\cellcolor[HTML]{F2F2F2}{ -1,50 }} & {\cellcolor[HTML]{F2F2F2}{ -0,92 }} & {\cellcolor[HTML]{F2F2F2}{ -0,48 }} & {\cellcolor[HTML]{F2F2F2}{ -0,09 }} & {\cellcolor[HTML]{EBEBEB}{  0,21 }} & {\cellcolor[HTML]{EBEBEB}{  0,57 }} & {\cellcolor[HTML]{EBEBEB}{  0,57 }} \\ 
Impulsion budgétaire primaire\textsuperscript{\textit{b}} &   0,32  &  -1,06  & {\cellcolor[HTML]{F2F2F2}{  0,00 }} & {\cellcolor[HTML]{F2F2F2}{ -1,15 }} & {\cellcolor[HTML]{F2F2F2}{ -0,58 }} & {\cellcolor[HTML]{F2F2F2}{ -0,44 }} & {\cellcolor[HTML]{F2F2F2}{ -0,38 }} & {\cellcolor[HTML]{EBEBEB}{ -0,30 }} & {\cellcolor[HTML]{EBEBEB}{ -0,00 }} & {\cellcolor[HTML]{EBEBEB}{—}} \\ 
   dont PO\textsuperscript{\textit{b}} &   0,36  &  -0,84  & {\cellcolor[HTML]{F2F2F2}{  0,00 }} & {\cellcolor[HTML]{F2F2F2}{ -0,57 }} & {\cellcolor[HTML]{F2F2F2}{ -0,29 }} & {\cellcolor[HTML]{F2F2F2}{ -0,22 }} & {\cellcolor[HTML]{F2F2F2}{ -0,19 }} & {\cellcolor[HTML]{EBEBEB}{ -0,15 }} & {\cellcolor[HTML]{EBEBEB}{ -0,00 }} & {\cellcolor[HTML]{EBEBEB}{—}} \\ 
   dont dépense\textsuperscript{\textit{b}} &  -0,03  &  -0,09  & {\cellcolor[HTML]{F2F2F2}{  0,00 }} & {\cellcolor[HTML]{F2F2F2}{ -0,57 }} & {\cellcolor[HTML]{F2F2F2}{ -0,29 }} & {\cellcolor[HTML]{F2F2F2}{ -0,22 }} & {\cellcolor[HTML]{F2F2F2}{ -0,19 }} & {\cellcolor[HTML]{EBEBEB}{ -0,15 }} & {\cellcolor[HTML]{EBEBEB}{ -0,00 }} & {\cellcolor[HTML]{EBEBEB}{—}} \\ 
Cumul de l'impulsion primaire\textsuperscript{\textit{b}} & — & — & {\cellcolor[HTML]{F2F2F2}{  0,00 }} & {\cellcolor[HTML]{F2F2F2}{ -1,15 }} & {\cellcolor[HTML]{F2F2F2}{ -1,72 }} & {\cellcolor[HTML]{F2F2F2}{ -2,16 }} & {\cellcolor[HTML]{F2F2F2}{ -2,54 }} & {\cellcolor[HTML]{EBEBEB}{ -2,84 }} & {\cellcolor[HTML]{EBEBEB}{ -3,20 }} & \textsuperscript{\textit{c}}{\cellcolor[HTML]{EBEBEB}{ -3,20 }} \\ 
{Taux souverain courant}\textsuperscript{\textit{f}} & {  2,97 } & {  3,35 } & {\cellcolor[HTML]{F2F2F2}{  3,53 }} & {\cellcolor[HTML]{F2F2F2}{  3,76 }} & {\cellcolor[HTML]{F2F2F2}{  3,52 }} & {\cellcolor[HTML]{F2F2F2}{  3,27 }} & {\cellcolor[HTML]{F2F2F2}{  3,28 }} & {\cellcolor[HTML]{EBEBEB}{  3,27 }} & {\cellcolor[HTML]{EBEBEB}{  3,29 }} & {\cellcolor[HTML]{EBEBEB}{—}} \\ 
Taux souverain pondéré\textsuperscript{\textit{g}} &   1,94  &   2,02  & {\cellcolor[HTML]{F2F2F2}{  2,11 }} & {\cellcolor[HTML]{F2F2F2}{  2,29 }} & {\cellcolor[HTML]{F2F2F2}{  2,43 }} & {\cellcolor[HTML]{F2F2F2}{  2,52 }} & {\cellcolor[HTML]{F2F2F2}{  2,61 }} & {\cellcolor[HTML]{EBEBEB}{  2,68 }} & {\cellcolor[HTML]{EBEBEB}{  3,24 }} & {\cellcolor[HTML]{EBEBEB}{—}} \\ 
Dette publique\textsuperscript{\textit{b}} & 113,23  & 115,60  & {\cellcolor[HTML]{F2F2F2}{117,32 }} & {\cellcolor[HTML]{F2F2F2}{119,85 }} & {\cellcolor[HTML]{F2F2F2}{121,18 }} & {\cellcolor[HTML]{F2F2F2}{121,92 }} & {\cellcolor[HTML]{F2F2F2}{122,22 }} & {\cellcolor[HTML]{EBEBEB}{122,26 }} & {\cellcolor[HTML]{EBEBEB}{120,10 }} & {\cellcolor[HTML]{EBEBEB}{122,26 }} \\ 
Charge d'intérêts\textsuperscript{\textit{b}} &   2,19  &   2,34  & {\cellcolor[HTML]{F2F2F2}{  2,37 }} & {\cellcolor[HTML]{F2F2F2}{  2,64 }} & {\cellcolor[HTML]{F2F2F2}{  2,84 }} & {\cellcolor[HTML]{F2F2F2}{  2,97 }} & {\cellcolor[HTML]{F2F2F2}{  3,09 }} & {\cellcolor[HTML]{EBEBEB}{  3,19 }} & {\cellcolor[HTML]{EBEBEB}{  3,78 }} & {\cellcolor[HTML]{EBEBEB}{  3,80 }} \\ 
\bottomrule
\end{tabular*}
\begin{minipage}{\linewidth}
\vspace{.05em}
\textsuperscript{\textit{a}}évolution en \% annuel, volume \textsuperscript{\textit{b}}\% du PIB \textsuperscript{\textit{c}}Minimum \textsuperscript{\textit{d}}déflateur du PIB, évolution en \% annuel \textsuperscript{\textit{e}}\% de la population active \textsuperscript{\textit{f}}10 ans, \% annuel \textsuperscript{\textit{g}}\% annuel\\
Données AMECO printemps 2026, première année de simulation : 2026\\
Simulation debtwatch v2.0 (github.com/OFCE/dwr), France, stochastique (256 tirages), valeurs moyennes\\
\\
\end{minipage}

}

\end{table}%

Cet ajustement limite la hausse de la dette publique à 122,2\% du PIB
(tableau~\ref{tbl-optimal1}), bien que l'intervalle de confiance soit
assez important (graphique~\ref{fig-optimal}), en particulier sur la
dette publique.

Notons que cet écart est construit sous des hypothèses particulières
(voir Timbeau 2025) et doit être interprété en fonction. En simplifiant
à l'extrême, il s'agit de la prise en compte de la variabilité
historique de la croissance, mais en supposant que la moyenne (de la
croissance) est stable et qu'il y a peu d'autocorrélation dans ce
processus (il est non persistant). Cette hypothèse induit un retour à la
moyenne assez sûr de l'écart de production et donc limite le
renforcement d'une déviation médianne de la trajectoire de dette.

L'ajustement nécessaire cumulé en 2031 est de 2,84~points de PIB soit
approximativement 85~milliards d'euros. Il est de l'ordre de
105~milliards d'euros à terme.

Le tableau~\ref{tbl-optimal2} décrit un monde légèrement différent. Au
lieu d'une impulsion nulle en 2026, on suppose une impulsion négative de
0,5~point de PIB. L'optimisation conduit à une impulsion plus lissée en
2027, ce qui est un résultat central dans le calcul d'une politique
budgétaire optimale\,: le \emph{backloading} n'est justifié qu'en cas
d'une forte préférence pour le présent, une préférence pour des
impulsions les moins variées possibles ou des mulitiplicateurs
budgétaires qui varient dans le cycle.

Le maximum de dette est un peu moins important. Le fait qu'avec un
ajustement plus précoce l'écart de taux avec des obligations de
référence puisse se réduire n'est pas pris en compte et accenturait le
constraste avec le scénario principal. Il existe cependant peu
d'éléments empiriques pour justifier une variation franche de cet écart,
et la nature \emph{forward-looking} des déterminants de l'écart de taux
incite à la prudence\,: pas grand chose ne sépare ces deux trajectoires
et donc elles ne sont pas plus rassurantes l'une que l'autre, tout au
plus la seconde produisant un ajustement plus précoce, elle pourrait
être associée à un taux souverain légèrement plus détendu.

\begin{table}

\caption{\label{tbl-optimal2}Ajustement optimal, -0,5\% d'impulsion en
2026}

\centering{

\fontsize{6.0pt}{7.0pt}\selectfont
\begin{tabular*}{\linewidth}{@{\extracolsep{\fill}}>{\raggedright\arraybackslash}p{\dimexpr 165.00pt -2\tabcolsep-1.5\arrayrulewidth}>{\raggedleft\arraybackslash}p{\dimexpr 37.50pt -2\tabcolsep-1.5\arrayrulewidth}>{\raggedleft\arraybackslash}p{\dimexpr 37.50pt -2\tabcolsep-1.5\arrayrulewidth}>{\raggedleft\arraybackslash}p{\dimexpr 37.50pt -2\tabcolsep-1.5\arrayrulewidth}>{\raggedleft\arraybackslash}p{\dimexpr 37.50pt -2\tabcolsep-1.5\arrayrulewidth}>{\raggedleft\arraybackslash}p{\dimexpr 37.50pt -2\tabcolsep-1.5\arrayrulewidth}>{\raggedleft\arraybackslash}p{\dimexpr 37.50pt -2\tabcolsep-1.5\arrayrulewidth}>{\raggedleft\arraybackslash}p{\dimexpr 37.50pt -2\tabcolsep-1.5\arrayrulewidth}>{\raggedleft\arraybackslash}p{\dimexpr 37.50pt -2\tabcolsep-1.5\arrayrulewidth}>{\raggedleft\arraybackslash}p{\dimexpr 37.50pt -2\tabcolsep-1.5\arrayrulewidth}>{\raggedleft\arraybackslash}p{\dimexpr 37.50pt -2\tabcolsep-1.5\arrayrulewidth}}
\toprule
 & 2024 & 2025 & 2026 & 2027 & 2028 & 2029 & 2030 & 2031 & LT (2050) & Max. 25-50 \\ 
\midrule\addlinespace[2.5pt]
Croissance du PIB\textsuperscript{\textit{a}} &   1,19  &   0,84  & {\cellcolor[HTML]{F2F2F2}{  0,62 }} & {\cellcolor[HTML]{F2F2F2}{  0,43 }} & {\cellcolor[HTML]{F2F2F2}{  1,06 }} & {\cellcolor[HTML]{F2F2F2}{  1,19 }} & {\cellcolor[HTML]{F2F2F2}{  1,21 }} & {\cellcolor[HTML]{EBEBEB}{  1,20 }} & {\cellcolor[HTML]{EBEBEB}{  1,01 }} & {\cellcolor[HTML]{EBEBEB}{—}} \\ 
Ecart de production\textsuperscript{\textit{b}} &  -0,28  &   0,21  & {\cellcolor[HTML]{F2F2F2}{ -0,34 }} & {\cellcolor[HTML]{F2F2F2}{ -1,07 }} & {\cellcolor[HTML]{F2F2F2}{ -1,16 }} & {\cellcolor[HTML]{F2F2F2}{ -1,02 }} & {\cellcolor[HTML]{F2F2F2}{ -0,84 }} & {\cellcolor[HTML]{EBEBEB}{ -0,67 }} & {\cellcolor[HTML]{EBEBEB}{  0,02 }} & \textsuperscript{\textit{c}}{\cellcolor[HTML]{EBEBEB}{ -1,16 }} \\ 
Croissance du PIB potentiel\textsuperscript{\textit{a}} & — & — & {\cellcolor[HTML]{F2F2F2}{  1,18 }} & {\cellcolor[HTML]{F2F2F2}{  1,16 }} & {\cellcolor[HTML]{F2F2F2}{  1,15 }} & {\cellcolor[HTML]{F2F2F2}{  1,04 }} & {\cellcolor[HTML]{F2F2F2}{  1,03 }} & {\cellcolor[HTML]{EBEBEB}{  1,03 }} & {\cellcolor[HTML]{EBEBEB}{  1,00 }} & {\cellcolor[HTML]{EBEBEB}{—}} \\ 
Inflation\textsuperscript{\textit{d}} &   2,09  &   1,18  & {\cellcolor[HTML]{F2F2F2}{  1,82 }} & {\cellcolor[HTML]{F2F2F2}{  1,76 }} & {\cellcolor[HTML]{F2F2F2}{  1,74 }} & {\cellcolor[HTML]{F2F2F2}{  1,60 }} & {\cellcolor[HTML]{F2F2F2}{  1,58 }} & {\cellcolor[HTML]{EBEBEB}{  1,55 }} & {\cellcolor[HTML]{EBEBEB}{  1,72 }} & {\cellcolor[HTML]{EBEBEB}{—}} \\ 
Taux de chômage\textsuperscript{\textit{e}} &   7,40  &   7,70  & {\cellcolor[HTML]{F2F2F2}{  8,03 }} & {\cellcolor[HTML]{F2F2F2}{  8,43 }} & {\cellcolor[HTML]{F2F2F2}{  8,41 }} & {\cellcolor[HTML]{F2F2F2}{  8,29 }} & {\cellcolor[HTML]{F2F2F2}{  8,10 }} & {\cellcolor[HTML]{EBEBEB}{  7,90 }} & {\cellcolor[HTML]{EBEBEB}{  6,98 }} & {\cellcolor[HTML]{EBEBEB}{  8,43 }} \\ 
{Solde public}\textsuperscript{\textit{b}} & { -5,93 } & { -5,09 } & {\cellcolor[HTML]{F2F2F2}{ -4,67 }} & {\cellcolor[HTML]{F2F2F2}{ -4,52 }} & {\cellcolor[HTML]{F2F2F2}{ -4,28 }} & {\cellcolor[HTML]{F2F2F2}{ -3,95 }} & {\cellcolor[HTML]{F2F2F2}{ -3,64 }} & {\cellcolor[HTML]{EBEBEB}{ -3,39 }} & {\cellcolor[HTML]{EBEBEB}{ -3,20 }} & {\cellcolor[HTML]{EBEBEB}{ -2,98 }} \\ 
Solde public primaire\textsuperscript{\textit{b}} &  -3,73  &  -2,75  & {\cellcolor[HTML]{F2F2F2}{ -2,29 }} & {\cellcolor[HTML]{F2F2F2}{ -1,88 }} & {\cellcolor[HTML]{F2F2F2}{ -1,46 }} & {\cellcolor[HTML]{F2F2F2}{ -1,00 }} & {\cellcolor[HTML]{F2F2F2}{ -0,57 }} & {\cellcolor[HTML]{EBEBEB}{ -0,22 }} & {\cellcolor[HTML]{EBEBEB}{  0,58 }} & {\cellcolor[HTML]{EBEBEB}{  0,59 }} \\ 
Solde public structurel primaire\textsuperscript{\textit{b}} &  -3,61  &  -2,85  & {\cellcolor[HTML]{F2F2F2}{ -2,13 }} & {\cellcolor[HTML]{F2F2F2}{ -1,38 }} & {\cellcolor[HTML]{F2F2F2}{ -0,91 }} & {\cellcolor[HTML]{F2F2F2}{ -0,52 }} & {\cellcolor[HTML]{F2F2F2}{ -0,17 }} & {\cellcolor[HTML]{EBEBEB}{  0,10 }} & {\cellcolor[HTML]{EBEBEB}{  0,57 }} & {\cellcolor[HTML]{EBEBEB}{  0,57 }} \\ 
Impulsion budgétaire primaire\textsuperscript{\textit{b}} &   0,32  &  -1,06  & {\cellcolor[HTML]{F2F2F2}{ -0,50 }} & {\cellcolor[HTML]{F2F2F2}{ -0,77 }} & {\cellcolor[HTML]{F2F2F2}{ -0,46 }} & {\cellcolor[HTML]{F2F2F2}{ -0,39 }} & {\cellcolor[HTML]{F2F2F2}{ -0,34 }} & {\cellcolor[HTML]{EBEBEB}{ -0,27 }} & {\cellcolor[HTML]{EBEBEB}{ -0,01 }} & {\cellcolor[HTML]{EBEBEB}{—}} \\ 
   dont PO\textsuperscript{\textit{b}} &   0,36  &  -0,84  & {\cellcolor[HTML]{F2F2F2}{ -0,25 }} & {\cellcolor[HTML]{F2F2F2}{ -0,38 }} & {\cellcolor[HTML]{F2F2F2}{ -0,23 }} & {\cellcolor[HTML]{F2F2F2}{ -0,19 }} & {\cellcolor[HTML]{F2F2F2}{ -0,17 }} & {\cellcolor[HTML]{EBEBEB}{ -0,14 }} & {\cellcolor[HTML]{EBEBEB}{ -0,00 }} & {\cellcolor[HTML]{EBEBEB}{—}} \\ 
   dont dépense\textsuperscript{\textit{b}} &  -0,03  &  -0,09  & {\cellcolor[HTML]{F2F2F2}{ -0,25 }} & {\cellcolor[HTML]{F2F2F2}{ -0,38 }} & {\cellcolor[HTML]{F2F2F2}{ -0,23 }} & {\cellcolor[HTML]{F2F2F2}{ -0,19 }} & {\cellcolor[HTML]{F2F2F2}{ -0,17 }} & {\cellcolor[HTML]{EBEBEB}{ -0,14 }} & {\cellcolor[HTML]{EBEBEB}{ -0,00 }} & {\cellcolor[HTML]{EBEBEB}{—}} \\ 
Cumul de l'impulsion primaire\textsuperscript{\textit{b}} & — & — & {\cellcolor[HTML]{F2F2F2}{ -0,50 }} & {\cellcolor[HTML]{F2F2F2}{ -1,27 }} & {\cellcolor[HTML]{F2F2F2}{ -1,73 }} & {\cellcolor[HTML]{F2F2F2}{ -2,12 }} & {\cellcolor[HTML]{F2F2F2}{ -2,46 }} & {\cellcolor[HTML]{EBEBEB}{ -2,73 }} & {\cellcolor[HTML]{EBEBEB}{ -3,20 }} & \textsuperscript{\textit{c}}{\cellcolor[HTML]{EBEBEB}{ -3,20 }} \\ 
{Taux souverain courant}\textsuperscript{\textit{f}} & {  2,97 } & {  3,35 } & {\cellcolor[HTML]{F2F2F2}{  3,53 }} & {\cellcolor[HTML]{F2F2F2}{  3,76 }} & {\cellcolor[HTML]{F2F2F2}{  3,52 }} & {\cellcolor[HTML]{F2F2F2}{  3,27 }} & {\cellcolor[HTML]{F2F2F2}{  3,28 }} & {\cellcolor[HTML]{EBEBEB}{  3,28 }} & {\cellcolor[HTML]{EBEBEB}{  3,29 }} & {\cellcolor[HTML]{EBEBEB}{—}} \\ 
Taux souverain pondéré\textsuperscript{\textit{g}} &   1,94  &   2,02  & {\cellcolor[HTML]{F2F2F2}{  2,11 }} & {\cellcolor[HTML]{F2F2F2}{  2,29 }} & {\cellcolor[HTML]{F2F2F2}{  2,43 }} & {\cellcolor[HTML]{F2F2F2}{  2,52 }} & {\cellcolor[HTML]{F2F2F2}{  2,61 }} & {\cellcolor[HTML]{EBEBEB}{  2,68 }} & {\cellcolor[HTML]{EBEBEB}{  3,24 }} & {\cellcolor[HTML]{EBEBEB}{—}} \\ 
Dette publique\textsuperscript{\textit{b}} & 113,23  & 115,60  & {\cellcolor[HTML]{F2F2F2}{117,51 }} & {\cellcolor[HTML]{F2F2F2}{119,51 }} & {\cellcolor[HTML]{F2F2F2}{120,53 }} & {\cellcolor[HTML]{F2F2F2}{121,20 }} & {\cellcolor[HTML]{F2F2F2}{121,54 }} & {\cellcolor[HTML]{EBEBEB}{121,67 }} & {\cellcolor[HTML]{EBEBEB}{120,10 }} & {\cellcolor[HTML]{EBEBEB}{121,67 }} \\ 
Charge d'intérêts\textsuperscript{\textit{b}} &   2,19  &   2,34  & {\cellcolor[HTML]{F2F2F2}{  2,38 }} & {\cellcolor[HTML]{F2F2F2}{  2,63 }} & {\cellcolor[HTML]{F2F2F2}{  2,82 }} & {\cellcolor[HTML]{F2F2F2}{  2,96 }} & {\cellcolor[HTML]{F2F2F2}{  3,07 }} & {\cellcolor[HTML]{EBEBEB}{  3,17 }} & {\cellcolor[HTML]{EBEBEB}{  3,78 }} & {\cellcolor[HTML]{EBEBEB}{  3,80 }} \\ 
\bottomrule
\end{tabular*}
\begin{minipage}{\linewidth}
\vspace{.05em}
\textsuperscript{\textit{a}}évolution en \% annuel, volume \textsuperscript{\textit{b}}\% du PIB \textsuperscript{\textit{c}}Minimum \textsuperscript{\textit{d}}déflateur du PIB, évolution en \% annuel \textsuperscript{\textit{e}}\% de la population active \textsuperscript{\textit{f}}10 ans, \% annuel \textsuperscript{\textit{g}}\% annuel\\
Données AMECO printemps 2026, première année de simulation : 2026\\
Simulation debtwatch v2.0 (github.com/OFCE/dwr), France, stochastique (256 tirages), valeurs moyennes\\
\\
\end{minipage}

}

\end{table}%

\begin{figure}[htb]

\caption{\label{fig-optimal}Ajustement optimal}

\centering{

\includegraphics[width=1\linewidth,height=\textheight,keepaspectratio]{note-ajustement-2026-XT_files/figure-pdf/fig-optimal-1.png}

}

\end{figure}%

\chapter{Annexes}\label{annexes}

\section*{Annexe 1\,: Principales
hypothèses}\label{annexe-1-principales-hypothuxe8ses}
\addcontentsline{toc}{section}{Annexe 1\,: Principales hypothèses}

\begin{enumerate}
\def\labelenumi{\arabic{enumi}.}
\item
  On a supposé une croissance potentielle de 1\% par an à terme. A moyen
  terme, la croissance potentielle est un peu supérieure du fait d'une
  baisse du chômage structurel vers 7\% de la population active.
\item
  L'inflation est ancrée à 1,75\% à long terme. En 2026, en suposant que
  le prix du pétrole est revenu à sa valeur d'avant la Guerre du Golfe
  III, l'inflation en 2026 en moyenne annuelle s'établirait juste au
  dessus de 1,8\% par an.
\item
  Le solde public et la dette publique sont calés sur la dernière
  notification de la dette publique française (mars 2026). Les autres
  données macroéconomique pour l'année 2025 sont celle de la base AMECO
  de mai 2026. Elles diffèrent des comptes nationaux de mai 2026 à la
  marge.
\item
  L'écart de production est calé sur l'analyse de la DGTrésor. AMECO
  donne un écart de production de -0.9\% du PIB potentiel. Cela
  ajouterait de l'ordre de 0,2\% de croissance à court terme dans les
  scénarios de projection.
\item
  La prime sur le taux souverain par rapport à la croissance potentielle
  nominale s'accroit encore en 2027 puis décroît lentement tout en
  restant au dessus de la moyenen observée les dernières années. Il y a
  des raisons de penser qu'il pourrait être significativement plus bas à
  l'horizon de 5~ans (de l'ordre d'un demi-point).
\item
  L'optimisation est une optimisation sous contrainte d'atteindre la
  cible de dette (120\%) à une date donnée (2037) et de s'y tenir
  ensuite. Les impulsions sont calculées en fonction en minimisant une
  fonction objectif qui minimise la perte quadratique d'écart de
  croissance et la variance des impulsions. Cette optmisation intègre
  les effets de retour des impulsions sur l'activité et donc les écarts
  de production.
\end{enumerate}

\section*{Annexe 2\,: Trajectoires au fil de
l'eau}\label{annexe-2-trajectoires-au-fil-de-leau}
\addcontentsline{toc}{section}{Annexe 2\,: Trajectoires au fil de l'eau}

La trajectoire au fil de l'eau est là à titre de référence. elle
consiste à dérouler à partir des conditions initiales la logique de la
dynamique du modèle sans impulsion budgétaire. Les propriétés du modèle
sont très simple, avec d'un côté le retour à l'équilibre de l'écart de
production et de l'autre le calcul de l'évolution de la dette publique
par l'accumulation des déficits structurels et de la charge d'intérêt.
L'inflation dépend, autour de la trajectoire, de l'écart de production
(que ce soit pas un arguement de type Phillips ou un ajustement par le
compte extérieur). La condition initiale en 2025 est celle d'un écart de
production peu ouvert (conforme à l'hypothèse DGTrésor). L'appréciation
de la position dans le cycle de la Commission Européenne est un peu plus
négative, ce qui justifierait un peu plus de croissance dans le court
terme.

La trajectorie au fil de l'eau est construite sur une impulsion
budgétaire nulle, c'est-à-dire que le taux de prélèvement obligatoire
est constant et que la dépense publique est constante en point de PIB
potentiel. Cette trajectoire est déjà un ajustement budgétaire si on
tient compte d'évolutions tendacielles ou votées comme la dépense
sociale, la loi de programation pluri-annuelle militaire ou les impôts
exceptionnels qui doivent s'éteindre en 2026 ou en 2027. Le fil de l'eau
choisi n'a donc rien de spontané mais illustre plutôt les dynamiques à
l'oeuvre.

De ce point de vue, la stabilité du déficit structurel ajouté à la
hausse de la charge d'intérêt conduit à un accroissement du déficit
public. A l'horizon de 2031, la hausse ainsi induite du déficit public
est de l'ordre d'un point de PIB, ce résultat dépendant des hypothèses
retenues sur les taux d'intérêt souverain courant et pondéré. Notons que
la modélisation du taux pondéré, avec un modèle auto-régressif, minore
potentiellement la hausse du taux pondéré à l'horizon de la décennie. A
plus long terme, cette déviation s'accentue, conduisant la dette
publique à 200\% du PIB en 2050. Ceci constitue indéniablement une
dynamique explosive mais relativement lente, en 25 années la dette
s'accroît de 85~points de PIB, soit 3,4~points par an.

\begin{table}

\caption{\label{tbl-fildle}Ajustement optimal, -0,5\% d'impulsion en
2026}

\centering{

\fontsize{6.0pt}{7.0pt}\selectfont
\begin{tabular*}{\linewidth}{@{\extracolsep{\fill}}>{\raggedright\arraybackslash}p{\dimexpr 165.00pt -2\tabcolsep-1.5\arrayrulewidth}>{\raggedleft\arraybackslash}p{\dimexpr 37.50pt -2\tabcolsep-1.5\arrayrulewidth}>{\raggedleft\arraybackslash}p{\dimexpr 37.50pt -2\tabcolsep-1.5\arrayrulewidth}>{\raggedleft\arraybackslash}p{\dimexpr 37.50pt -2\tabcolsep-1.5\arrayrulewidth}>{\raggedleft\arraybackslash}p{\dimexpr 37.50pt -2\tabcolsep-1.5\arrayrulewidth}>{\raggedleft\arraybackslash}p{\dimexpr 37.50pt -2\tabcolsep-1.5\arrayrulewidth}>{\raggedleft\arraybackslash}p{\dimexpr 37.50pt -2\tabcolsep-1.5\arrayrulewidth}>{\raggedleft\arraybackslash}p{\dimexpr 37.50pt -2\tabcolsep-1.5\arrayrulewidth}>{\raggedleft\arraybackslash}p{\dimexpr 37.50pt -2\tabcolsep-1.5\arrayrulewidth}>{\raggedleft\arraybackslash}p{\dimexpr 37.50pt -2\tabcolsep-1.5\arrayrulewidth}>{\raggedleft\arraybackslash}p{\dimexpr 37.50pt -2\tabcolsep-1.5\arrayrulewidth}}
\toprule
 & 2024 & 2025 & 2026 & 2027 & 2028 & 2029 & 2030 & 2031 & LT (2050) & Max. 25-50 \\ 
\midrule\addlinespace[2.5pt]
Croissance du PIB\textsuperscript{\textit{a}} &   1,19  &   0,84  & {\cellcolor[HTML]{F2F2F2}{  1,05 }} & {\cellcolor[HTML]{F2F2F2}{  1,02 }} & {\cellcolor[HTML]{F2F2F2}{  1,19 }} & {\cellcolor[HTML]{F2F2F2}{  1,10 }} & {\cellcolor[HTML]{F2F2F2}{  1,06 }} & {\cellcolor[HTML]{EBEBEB}{  1,04 }} & {\cellcolor[HTML]{EBEBEB}{  1,00 }} & {\cellcolor[HTML]{EBEBEB}{—}} \\ 
Ecart de production\textsuperscript{\textit{b}} &  -0,28  &   0,21  & {\cellcolor[HTML]{F2F2F2}{  0,08 }} & {\cellcolor[HTML]{F2F2F2}{ -0,06 }} & {\cellcolor[HTML]{F2F2F2}{ -0,02 }} & {\cellcolor[HTML]{F2F2F2}{  0,04 }} & {\cellcolor[HTML]{F2F2F2}{  0,06 }} & {\cellcolor[HTML]{EBEBEB}{  0,07 }} & {\cellcolor[HTML]{EBEBEB}{  0,10 }} & \textsuperscript{\textit{c}}{\cellcolor[HTML]{EBEBEB}{ -0,06 }} \\ 
Croissance du PIB potentiel\textsuperscript{\textit{a}} & — & — & {\cellcolor[HTML]{F2F2F2}{  1,18 }} & {\cellcolor[HTML]{F2F2F2}{  1,16 }} & {\cellcolor[HTML]{F2F2F2}{  1,15 }} & {\cellcolor[HTML]{F2F2F2}{  1,04 }} & {\cellcolor[HTML]{F2F2F2}{  1,03 }} & {\cellcolor[HTML]{EBEBEB}{  1,03 }} & {\cellcolor[HTML]{EBEBEB}{  1,00 }} & {\cellcolor[HTML]{EBEBEB}{—}} \\ 
Inflation\textsuperscript{\textit{d}} &   2,09  &   1,18  & {\cellcolor[HTML]{F2F2F2}{  1,82 }} & {\cellcolor[HTML]{F2F2F2}{  1,78 }} & {\cellcolor[HTML]{F2F2F2}{  1,82 }} & {\cellcolor[HTML]{F2F2F2}{  1,74 }} & {\cellcolor[HTML]{F2F2F2}{  1,73 }} & {\cellcolor[HTML]{EBEBEB}{  1,69 }} & {\cellcolor[HTML]{EBEBEB}{  1,73 }} & {\cellcolor[HTML]{EBEBEB}{—}} \\ 
Taux de chômage\textsuperscript{\textit{e}} &   7,40  &   7,70  & {\cellcolor[HTML]{F2F2F2}{  7,78 }} & {\cellcolor[HTML]{F2F2F2}{  7,70 }} & {\cellcolor[HTML]{F2F2F2}{  7,39 }} & {\cellcolor[HTML]{F2F2F2}{  7,23 }} & {\cellcolor[HTML]{F2F2F2}{  7,14 }} & {\cellcolor[HTML]{EBEBEB}{  7,08 }} & {\cellcolor[HTML]{EBEBEB}{  6,91 }} & {\cellcolor[HTML]{EBEBEB}{  7,78 }} \\ 
{Solde public}\textsuperscript{\textit{b}} & { -5,93 } & { -5,09 } & {\cellcolor[HTML]{F2F2F2}{ -4,95 }} & {\cellcolor[HTML]{F2F2F2}{ -5,27 }} & {\cellcolor[HTML]{F2F2F2}{ -5,45 }} & {\cellcolor[HTML]{F2F2F2}{ -5,59 }} & {\cellcolor[HTML]{F2F2F2}{ -5,73 }} & {\cellcolor[HTML]{EBEBEB}{ -5,88 }} & {\cellcolor[HTML]{EBEBEB}{ -8,47 }} & {\cellcolor[HTML]{EBEBEB}{ -4,95 }} \\ 
Solde public primaire\textsuperscript{\textit{b}} &  -3,73  &  -2,75  & {\cellcolor[HTML]{F2F2F2}{ -2,58 }} & {\cellcolor[HTML]{F2F2F2}{ -2,65 }} & {\cellcolor[HTML]{F2F2F2}{ -2,63 }} & {\cellcolor[HTML]{F2F2F2}{ -2,61 }} & {\cellcolor[HTML]{F2F2F2}{ -2,59 }} & {\cellcolor[HTML]{EBEBEB}{ -2,59 }} & {\cellcolor[HTML]{EBEBEB}{ -2,58 }} & {\cellcolor[HTML]{EBEBEB}{ -2,57 }} \\ 
Solde public structurel primaire\textsuperscript{\textit{b}} &  -3,61  &  -2,85  & {\cellcolor[HTML]{F2F2F2}{ -2,62 }} & {\cellcolor[HTML]{F2F2F2}{ -2,63 }} & {\cellcolor[HTML]{F2F2F2}{ -2,62 }} & {\cellcolor[HTML]{F2F2F2}{ -2,62 }} & {\cellcolor[HTML]{F2F2F2}{ -2,62 }} & {\cellcolor[HTML]{EBEBEB}{ -2,62 }} & {\cellcolor[HTML]{EBEBEB}{ -2,62 }} & {\cellcolor[HTML]{EBEBEB}{ -2,62 }} \\ 
Impulsion budgétaire primaire\textsuperscript{\textit{b}} &   0,32  &  -1,06  & {\cellcolor[HTML]{F2F2F2}{  0,00 }} & {\cellcolor[HTML]{F2F2F2}{  0,00 }} & {\cellcolor[HTML]{F2F2F2}{  0,00 }} & {\cellcolor[HTML]{F2F2F2}{  0,00 }} & {\cellcolor[HTML]{F2F2F2}{  0,00 }} & {\cellcolor[HTML]{EBEBEB}{  0,00 }} & {\cellcolor[HTML]{EBEBEB}{  0,00 }} & {\cellcolor[HTML]{EBEBEB}{—}} \\ 
   dont PO\textsuperscript{\textit{b}} &   0,36  &  -0,84  & {\cellcolor[HTML]{F2F2F2}{  0,00 }} & {\cellcolor[HTML]{F2F2F2}{  0,00 }} & {\cellcolor[HTML]{F2F2F2}{  0,00 }} & {\cellcolor[HTML]{F2F2F2}{  0,00 }} & {\cellcolor[HTML]{F2F2F2}{  0,00 }} & {\cellcolor[HTML]{EBEBEB}{  0,00 }} & {\cellcolor[HTML]{EBEBEB}{  0,00 }} & {\cellcolor[HTML]{EBEBEB}{—}} \\ 
   dont dépense\textsuperscript{\textit{b}} &  -0,03  &  -0,09  & {\cellcolor[HTML]{F2F2F2}{  0,00 }} & {\cellcolor[HTML]{F2F2F2}{  0,00 }} & {\cellcolor[HTML]{F2F2F2}{  0,00 }} & {\cellcolor[HTML]{F2F2F2}{  0,00 }} & {\cellcolor[HTML]{F2F2F2}{  0,00 }} & {\cellcolor[HTML]{EBEBEB}{  0,00 }} & {\cellcolor[HTML]{EBEBEB}{  0,00 }} & {\cellcolor[HTML]{EBEBEB}{—}} \\ 
Cumul de l'impulsion primaire\textsuperscript{\textit{b}} & — & — & {\cellcolor[HTML]{F2F2F2}{  0,00 }} & {\cellcolor[HTML]{F2F2F2}{  0,00 }} & {\cellcolor[HTML]{F2F2F2}{  0,00 }} & {\cellcolor[HTML]{F2F2F2}{  0,00 }} & {\cellcolor[HTML]{F2F2F2}{  0,00 }} & {\cellcolor[HTML]{EBEBEB}{  0,00 }} & {\cellcolor[HTML]{EBEBEB}{  0,00 }} & \textsuperscript{\textit{c}}{\cellcolor[HTML]{EBEBEB}{  0,00 }} \\ 
{Taux souverain courant}\textsuperscript{\textit{f}} & {  2,97 } & {  3,35 } & {\cellcolor[HTML]{F2F2F2}{  3,53 }} & {\cellcolor[HTML]{F2F2F2}{  3,76 }} & {\cellcolor[HTML]{F2F2F2}{  3,54 }} & {\cellcolor[HTML]{F2F2F2}{  3,30 }} & {\cellcolor[HTML]{F2F2F2}{  3,31 }} & {\cellcolor[HTML]{EBEBEB}{  3,30 }} & {\cellcolor[HTML]{EBEBEB}{  3,30 }} & {\cellcolor[HTML]{EBEBEB}{—}} \\ 
Taux souverain pondéré\textsuperscript{\textit{g}} &   1,94  &   2,02  & {\cellcolor[HTML]{F2F2F2}{  2,11 }} & {\cellcolor[HTML]{F2F2F2}{  2,29 }} & {\cellcolor[HTML]{F2F2F2}{  2,43 }} & {\cellcolor[HTML]{F2F2F2}{  2,53 }} & {\cellcolor[HTML]{F2F2F2}{  2,61 }} & {\cellcolor[HTML]{EBEBEB}{  2,69 }} & {\cellcolor[HTML]{EBEBEB}{  3,24 }} & {\cellcolor[HTML]{EBEBEB}{—}} \\ 
Dette publique\textsuperscript{\textit{b}} & 113,23  & 115,60  & {\cellcolor[HTML]{F2F2F2}{117,32 }} & {\cellcolor[HTML]{F2F2F2}{119,37 }} & {\cellcolor[HTML]{F2F2F2}{121,30 }} & {\cellcolor[HTML]{F2F2F2}{123,53 }} & {\cellcolor[HTML]{F2F2F2}{125,91 }} & {\cellcolor[HTML]{EBEBEB}{128,44 }} & {\cellcolor[HTML]{EBEBEB}{190,46 }} & {\cellcolor[HTML]{EBEBEB}{210,47 }} \\ 
Charge d'intérêts\textsuperscript{\textit{b}} &   2,19  &   2,34  & {\cellcolor[HTML]{F2F2F2}{  2,37 }} & {\cellcolor[HTML]{F2F2F2}{  2,61 }} & {\cellcolor[HTML]{F2F2F2}{  2,82 }} & {\cellcolor[HTML]{F2F2F2}{  2,98 }} & {\cellcolor[HTML]{F2F2F2}{  3,14 }} & {\cellcolor[HTML]{EBEBEB}{  3,30 }} & {\cellcolor[HTML]{EBEBEB}{  5,90 }} & {\cellcolor[HTML]{EBEBEB}{  6,19 }} \\ 
\bottomrule
\end{tabular*}
\begin{minipage}{\linewidth}
\vspace{.05em}
\textsuperscript{\textit{a}}évolution en \% annuel, volume \textsuperscript{\textit{b}}\% du PIB \textsuperscript{\textit{c}}Minimum \textsuperscript{\textit{d}}déflateur du PIB, évolution en \% annuel \textsuperscript{\textit{e}}\% de la population active \textsuperscript{\textit{f}}10 ans, \% annuel \textsuperscript{\textit{g}}\% annuel\\
Données AMECO printemps 2026, première année de simulation : 2026\\
Simulation debtwatch v2.0 (github.com/OFCE/dwr), France, stochastique (256 tirages), valeurs moyennes\\
\\
\end{minipage}

}

\end{table}%

\begin{figure}[htb]

\caption{\label{fig-fildle}Trajectoire au fil de l'eau}

\centering{

\includegraphics[width=1\linewidth,height=\textheight,keepaspectratio]{note-ajustement-2026-XT_files/figure-pdf/fig-fildle-1.png}

}

\end{figure}%

\section*{Annexe 2\,: Principes de la
modélisation}\label{annexe-2-principes-de-la-moduxe9lisation}
\addcontentsline{toc}{section}{Annexe 2\,: Principes de la modélisation}

L'équation de retour à l'équilibre de l'écart de production est écrite
comme un modèle à correction d'erreur (MCE). Les impulsions jouent sur
l'écart de production. Le multiplicateur instantané est défini par le
paramètre \(og_{dep}\) ou \(og_{po}\). Le multiplicateur complet (ou de
long terme) dépend des autres paramètres du MCE, \(og_{lag}\) et
\(og_{mce}\). Les valeurs utilisées pour les simulations sont\,:
\(og_{dep}=0.8\), \(og_{po}=0.5\), \(og_{mce}=0.3\), \(og_{lag}=0.25\).

\begin{equation}\protect\phantomsection\label{eq-og}{
\Delta og = og_{lag} \times \Delta og_{-1} -og_{mce} \times og + og_{dep} \times ib_{dep} + og_{po} \times ib_{po} + \varepsilon
}\end{equation}

Un bruit est ajouté, avec une autocorrélation. \(ogn_\sigma=0.015\) et
\(ogn_{ar}=0.9\).

\[
\Delta \varepsilon = ogn_{ar} \times \varepsilon_{-1} + \epsilon_{og} (\sim \mathcal{N}(0,ogn_{\sigma}^2))
\]

Le taux de croissance du prix du PIB (appelé inflation, on ne distingue
pas prix de consommation ou autres), suit également un MCE. Le long
terme est défini par la cible d'inflation et un effet de Phillips. Les
paramètres sont \(inf_{lag}=0\), \(inf_{mce}=0.68\),
\(inf_{tcho}=0.15\).

\[
\Delta\dot{p} = inf_{lag} \times d\dot p_{-1} - inf_{mce} \times (\dot p - inf_{star} - inf_{tcho} \times (u - nairu)) +\eta
\]

Comme pour l'écart de production, on ajoute une perturbation aléatoire.
\(inf_{\sigma}=0.01\), \(infn_{ar}=0.0\).

\[
\Delta \eta = infn_{ar} \times \eta_{-1} + \epsilon_{inf} (\sim \mathcal{N}(0,inf_{sigma}^2))
\]

La dynamique de la dette suit une équation habituelle. \(t_{aut}\)
assure l'équilibre comptable et prend en compte les recettes non
fiscales.

\[
\Delta dette = -s_p + int - t_{aut} \times vpib + one_{off}
\]

Le taux de marché est défini comme la somme de l'inflation, de la
croissance potentielle, d'une prime de risque \(ec\) et d'un terme qui
dépend de l'écart de la dette à une dette de référence \(d^{**}\). Dans
les simulations \(r_{dette}\) est nul.

\[
r_{inst} = g_{pot} + \dot{p} + ec + r_{dette} \times (d - d^{**})
\]

La prime de risque suit un MCE, \(ec_{mce}=0.5\).

\[
\Delta ec = -ec_{mce} \times (ec_{-1} - ec^*)
\]

Le taux moyen pondéré de la dette publique ou taux apparent est la
moyenne entre le taux moyen pondéré antérieur et le taux courant, comme
approximation de tranches de dettes de maturité égale et de différentes
dates d'émission. \(r_{mat}=9\).

\[
r_{app} = 1/r_{mat} \times r_{inst, -1} + (1-1/r_{mat}) \times r_{app, -1}
\]

La règle linéaire qui définit les impulsions est (\(tpo_{s^*}=0.25\),
\(tpo_{og}=-0.25\), \(tpo_{d^*}=-0.1\):

\[
ib_{lr} = tpo_{s^*} \times (spp - s^*) + tpo_{d^*} \times (d- d^*) + tpo_{og}\times og
\]

Le chômage suit un MCE dont la cible est le nairu et qui est perturbé
par l'écart de production, le coefficient d'Okun à court terme est
définit par \(tcho_{okun}=0.5\) et unitaire à long terme (i.e.~en cas de
choc persistant sur l'écart de production). \(tcho_{mce}=0.27\).

\[
\Delta tcho = - tcho_{okun} \times dog - tcho_{mce} \times (tcho+og-nairu)
\]

Le nairu s'ajuste progressivement à sa cible, afin de lisser les
évolutions. \(nairu_{mce}=0.2\)

\[
\Delta nairu = - nairu_{mce} \times (nairu_{-1} - nairu^*)
\]

Une politique budgétaire optimale est calculée\,: les impulsions
budgétaires successives (disons de 2025 à 2055) sont déterminées de
façon à ce que la dette publique soit égale à la cible de dette à partir
d'une date donnée (dans le cas central 2045, 20~ans après 2025). Ce
problème admet une infinité de solutions, puisque il suffit que le
déficit primaire soit nul à partir de 2045 et que la dette en 2045 soit
égale à la cible. On choisit parmi toutes ces solutions celle qui
minimise l'écart quadratique d'activité ou des impulsions. Formellement,
la fonction de perte s'écrit\,:

\begin{equation}\protect\phantomsection\label{eq-perte-quad}{
\mathcal{L} = \lambda_{og}\sum_{t=2025}^{\infty} \beta^t og_t^2 + \lambda_{ib} \sum_{t=2025}^{\infty} \beta^t ib_t^2 + \lambda_{max} \max_{t\geq 2025}(-og_t)
}\end{equation}

\(\lambda_{og}\), \(\lambda_{ib}\) et \(\lambda_{max}\) permettent de
donner un poids relatif aux trois termes de cette expression. \(\beta\)
est un facteur d'actualisation. Cette fonction de perte est permet
l'utilisation d'algorithmes numériques standards de minimisation (en
l'occurence COBYLA tel qu'implémenté dans Johnson (2008)).

Le premier terme s'interprète la minimisation de la «\,variance\,» des
écarts de production, c'est-à-dire une trajectoire qui s'écarte le moins
possible de 0 et qui recherche des trajectoires où l'écart est le plus
constant possible. On minimise ainsi l'impact de la politique
restrictive. La partie en \(ib_t^2\) est l'équivalent, mais pour les
instruments. Elle permet de représenter un monde où il est difficile de
modifier le budget. Cette inertie, si elle n'est pas rationnelle, est
réaliste. Le mélange entre les deux fonctions de perte conduit à
associer les deux arguments. Le facteur d'actualisation consiste à
préférer des conséquences lointaines plutôt qu'immédiates et constitue
une façon simple de justifier la \emph{backloading} par rapport au
\emph{frontloading}.

Enfin, le troisième terme conduit à prrendre en compte le creux d'écart
de production. Cette fonction de perte conduit, comme on le verra plus
loin, à produire des trajectoires dont la partie \(og_t\) négative est
la plus limitée possible. Ce terme de la fonction de perte induit un
\emph{backloading} assez fort, par construction. Pour faciliter le
calcul numérique, on utilise une approximation continue de la fonction
\(max\)\footnote{par une fonction
  \(max(x_1,x_2) \approx (x_1^\sigma + x_2^\sigma)^{1/\sigma}\) avec
  \(sigma\) assez grand (en pratique on prendra 8).}.

\begin{equation}\protect\phantomsection\label{eq-perte-min}{
\mathcal{L} = max(-og_t)
}\end{equation}

Le programme que l'on résout est alors donné par
équation~\ref{eq-prog-min}, \(\mathcal{S}\) étant la solution du système
dynamique qui permet de calculer la trajectoire de la dette à partir des
impulsions\,:

\begin{equation}\protect\phantomsection\label{eq-prog-min}{
\min_{ib_{t}}\mathcal{L} \\
\text{sous contrainte : pour } t\geq T \text{  } d_t = d^* \\
\text{avec } d_t = \mathcal{S}(ib_t)
}\end{equation}

\protect\phantomsection\label{refs}
\begin{CSLReferences}{1}{1}
\bibitem[\citeproctext]{ref-nloptr}
Johnson, Steven G. 2008. \emph{The NLopt nonlinear-optimization
package}. \url{https://github.com/stevengj/nlopt}.

\bibitem[\citeproctext]{ref-timbeau2025}
Timbeau, Xavier. 2025. {«~Trajectoires d'ajustement des finances
publiques de la France, une analyse avec debtwatch~»}. \emph{document de
travail de l'OFCE}, nᵒ 14 (juillet).
\url{https://ofce.github.io/debtwatch}.

\end{CSLReferences}




\end{document}
